% =====================================================================
% P5 merge snippet — M04: researchable questions with acceptance tests
% Source: manuscript_v2, Sec. 16 (tab:research-questions, R1--R10), written 2026-11-01
% as a consolidation of that paper's per-topic research questions.
% Adapted: the "stated in" column now points at THIS paper's sections (sec:open,
% sec:synthesis, sec:interfaces, sec:size, sec:methods) and appendices.
% Nothing new is claimed: each row restates a gap this paper already states in prose.
% NEW LABELS: tab:research-questions, sec:research_questions
% Suggested placement: end of \cref{sec:open}, after the prose agenda.
% =====================================================================

\subsection{Researchable questions, and what would settle each}\label{sec:research_questions}

The agenda above states the gaps. A gap becomes usable when the evidence that would close it is
stated too, because that is what lets a reader decide whether a proposed study is decisive or merely
related. The table makes the test explicit; the final column names the section where the gap is first
stated in this paper.

\begin{table}[H]\centering\small
\caption{Researchable questions. Each row restates a gap stated in \cref{sec:open}; the third column
names the evidence that would move it from open to closed. No new claim is introduced.}
\label{tab:research-questions}
\begin{tabularx}{\textwidth}{@{}p{0.5cm}Y Y p{1.75cm}@{}}
\toprule
\textbf{\#} & \textbf{Question} & \textbf{What would settle it} & \textbf{Stated in} \\
\midrule
R1 & What is the total (not availability) energy flux for each generalized family, and how does it
  differ from the flux used for interface partition? & A time-domain derivation per family under an
  explicit free-energy convention, external-field power terms and boundary conditions, checked against
  the primary formulation --- and a statement of the first-order reduction that recovers the flux used
  for the energy ratios. & \cref{sec:foundations}, \cref{sec:interfaces} \\
R2 & Is energy conserved at an interface in generalized and nonlocal media, and in what sense? & An
  interface balance derived from the selected model's flux, with the admissible jump conditions, tested
  against a benchmark whose reference solution does not depend on the model under test. &
  \cref{sec:interfaces} \\
R3 & Which benchmark is the right reference for a generalized theory? & A specification whose reference
  solution is independent --- analytic limit, symbolic determinant, or an independently implemented
  discretization --- with a stated tolerance. & \cref{app:registry} \\
R4 & Do the predicted generalized modes exist in a real material, with the predicted attenuation? & An
  experimental measurement of the predicted observables under conditions inside the model's stated
  validity domain. & \cref{sec:experiments} \\
R5 & How should energy flux and energy partition be normalized across families so that comparisons mean
  something? & A normalization in terms of physical observables rather than model parameters. &
  \cref{sec:synthesis} \\
R6 & How are nonlocal interface conditions and a nonlocal divergence theorem formulated? & A variational
  statement of the nonlocal boundary-value problem from which the interface conditions follow rather than
  are assumed. & \cref{sec:size} \\
R7 & Are phase-lag, memory-kernel and anisotropy parameters identifiable from measurable data? & The
  number and type of independent measurements needed to separate the parameters, established by analysis
  of the sensitivity structure or by a designed experiment; the answer is model-specific. &
  \cref{sec:experiments} \\
R8 & Do the thermodynamic restrictions reduce the admissible parameter set under anisotropic symmetry as
  they do in the isotropic case? & The anisotropic admissibility analysis carried through for each family,
  with the coordinate-free form of the fractional and memory operators under anisotropic differentiation.
  & \cref{sec:foundations} \\
R9 & Which comparative statements survive when each family's own conventions are respected? & A
  model-by-model full-text check of constitutive conventions, free energies and external-field power terms
  before any flux or mode count is compared. & \cref{sec:synthesis} \\
R10 & Can the benchmarks be reproduced from published inputs, and at what convergence rate? & Shared
  implementations, published inputs and reported convergence rates for every benchmark. &
  \cref{sec:methods} \\
\bottomrule
\end{tabularx}
\end{table}

Two of these questions are the same question at different depths: R1 and R9. A family-level flux
comparison cannot be trusted until the model-specific derivations behind it exist, which is why the
register in \cref{sec:eqregister} marks the affected entries open instead of filling them.
